\chapter{A first model}\label{chap:start}

This chapter builds a complete model on the synthetic SAM that comes with the
package, removes its import tariffs under two closures, and reads the results.
Every step is one command, and each is described in Part~\ref{part:ref}.

\section{The SAM}

\manref{cge_synthetic} loads the standard synthetic SAM as data: one
observation per account, a string variable with the account labels and one
numeric variable per account. It has three activities and three commodities
(agriculture, manufactures and services), labour and capital, one household,
the government with direct, sales and import taxes, savings--investment and the
rest of the world. \manref{cge_sam} reads it into a new model:
\stlog{gs_sam}
Cell (r, c) of the SAM is a payment from the column account c to the row
account r. The listing shows the make matrix (activities selling their output
to commodities) and the intermediate inputs (commodities bought by
activities). The SAM balances: every account's receipts equal its payments.

\section{Roles}

\manref{cge_map} assigns each account to a role of the model and checks that
every non-zero cell is a flow the model represents
(section~\ref{sec:flows}). Activities and commodities are paired by position,
each activity with its main commodity:
\stlog{gs_map}

\section{Calibration and checks}

The model's elasticities have no defaults. Here value added has an elasticity
of substitution of 0.8 between labour and capital, and the Armington and CET
elasticities are 2 for every commodity; \manref{cge_elasticities} gives
published values with their sources. \manref{cge_calibrate} computes every
parameter and the base value of every variable from the SAM and the
elasticities (chapter~\ref{chap:calibration}), sets the default closure, and
checks that the base solves every equation; \manref{cge_check} runs the
validation checks of chapter~\ref{chap:checks}:
\stlog{gs_calibrate}
The model has 93 equations. The largest scaled residual at the base is at the
precision of double arithmetic; doubling every fixed nominal value doubles
every price and leaves every quantity unchanged; and the solver returns to the
base from starting values perturbed by up to 10 per cent.

\section{A tariff removal}

\manref{cge_shock} sets every import tariff to zero, and \manref{cge_solve}
solves the model and stores the solution as the simulation \cmd{notariff}:
\stlog{gs_solve}
\manref{cge_results} summarises it, as percentage changes from the base:
\stlog{gs_results}
The default closure has a flexible exchange rate with fixed foreign savings,
fixed tax rates with flexible government savings, savings-driven investment,
the CPI as numeraire, and both factors fully employed and mobile
(chapter~\ref{chap:closures}). Under it, the exchange rate depreciates by 4.8
per cent, and imports and exports both rise. Government savings, 3 in the
base, fall by 271 per cent, to a deficit of about 5: the government loses its
tariff revenue. Investment, which adjusts to available savings, falls by 11.9
per cent. The household's equivalent variation is 2.67 per cent of its base
consumption.

\section{The same shock under another closure}

The closure matters as much as the shock. With government savings fixed and
direct tax rates adjusting to keep them, the tariff revenue is replaced by a
higher direct tax:
\stlog{gs_closure}
\manref{cge_report} shows the two simulations side by side:
\stlog{gs_report}
Direct tax revenue now rises by 34.9 per cent, government savings are
unchanged, investment rises by 2.7 per cent, and real consumption falls by 0.4
per cent; the equivalent variation is $-0.42$ per cent of consumption. Under
both closures real GDP changes little and the exchange rate depreciates by
about 4.8 per cent. The two simulations differ in who bears the loss of tariff
revenue: investment under the first closure, households under the second.

\section{Next steps}

\manref{cge_report} with \cmd{saving()} writes every table to Excel, and
\manref{cge_save} stores the model with its simulations. A real application
starts from a country's SAM, prepared with the SAM tools of
chapter~\ref{chap:samtools}, and states the source of every elasticity, the
closure and the sensitivity of the main results to the elasticities.
Chapter~\ref{chap:example} takes an unbalanced SAM through every step.
